\statement{C.40}{Proved}{The correctly normalized third cusp holds in every symmetric-four quotient}
Let \(\alpha=(-1,f)\) be an involution of the full
\(\widehat{\mathrm{GT}}_0\), and let \(u\in F'\) be its
canonical full R1 word. In the universal full dihedral quotient
\[
 D=\widehat{\mathbb Z}(t)\rtimes\langle a\rangle,\qquad
 a^2=1,\quad ata=t^{-1},\qquad x\longmapsto a,\quad y\longmapsto at,
\]
write \(u\longmapsto t^\kappa\), and set
\[
                         d=2\kappa+1.
 \tag{C.40.1}
\]
By A.45, \(d\) is a full unit. Since \(D'=\overline{\langle
t^2\rangle}\) and \(u\in F'\), one has \(\kappa\in2\widehat{\mathbb Z}\)
and \(d\equiv1\pmod4\). This is the global-sign normalization
of the possible multipliers \(\{d,-d\}\).

For every continuous homomorphism \(q:F\to S_4\),
\[
 q\bigl(x^d u^{-1}y^d u\bigr)\sim_{S_4}q\bigl((xy)^d\bigr).
 \tag{C.40.2}
\]
Moreover, if \(J_4\) is the intersection of all epimorphism
kernels \(F\to S_4\), both the normalized substitution
\(\theta_d=(d,u)\) and \(\psi=(1,u)\) induce automorphisms
of the actual joint finite quotient \(F/J_4\).
The result asserts no full or relative third-cusp identity.

\begin{proof}
\emph{The common symmetric-three values are forced.}\par
In \(S_3\), choose a reflection \(a\) and a rotation \(t\)
with \(ata=t^{-1}\). Write
\[
       u(a,at)=t^k,\qquad u(t,a)=t^\ell,\qquad
                  k,\ell\in\mathbb F_3.
\]
Here \(k\) is the reduction of \(\kappa\) in (C.40.1).
A.44 gives \(k\in\{0,2\}\), and the induced negative action
on this two-reflection quotient is identity. Its normalized
conjugating word has
\[
                          f(a,at)=1.
\]
On the mixed rotation/reflection quotient the negative action
inverts \(C_3\), irrespective of its reflection conjugator.
The full equation \(u=f\alpha(u)\) therefore gives
\[
 f(t,a)=t^{2\ell},\qquad
 u(a,t)=t^{-\ell},\qquad f(a,t)=t^{-2\ell}.
 \tag{C.40.3}
\]
The latter two formulas use full duality of \(u\) and \(f\),
proved for the canonical \(u\) in C.35. Every generating
ordered pair in \(S_3\) is an inner transform of exactly one
of the three displayed pair types. Thus \(k,\ell\), with
\[
              (k,\ell)\in\{0,2\}\times\mathbb F_3,
 \tag{C.40.4}
\]
determine all the common \(S_3\) values needed below.
In particular \(d\) reduces to \(1\pmod{12}\) when \(k=0\),
and to \(5\pmod{12}\) when \(k=2\).

\emph{An exact binary matrix calculation.}\par
Put \(G=\operatorname{GL}_2(\mathbb F_3)\),
\(L=\operatorname{SL}_2(\mathbb F_3)\), and
\(Q=G/\{\pm1\}\simeq S_4\). The last identification is
the faithful projective action on the four points of
\(\mathbb P^1(\mathbb F_3)\). Conjugation on the three
nonidentity elements of the projective normal Klein group
gives \(Q\to S_3\). Both \(f\) and \(u\) have values in
\(L\). A commutator word is unchanged when either input
is multiplied by the central sign: its two continuous
abelian exponents are zero. Consequently its values on
all generating pairs of \(G\) are determined by its values
on the following nine projective-pair representatives,
with transport by conjugation in \(G\).

For compactness the four entries of a matrix are listed
row by row. All entries in this table lie in \(\mathbb F_3\).
\[
\begin{array}{c|c|c|c}
j&A_j&B_j&(|\bar A_j|,|\bar B_j|,|\bar A_j\bar B_j|)\\ \hline
0&(0,1,1,0)&(0,1,1,1)&(2,4,3)\\
1&(0,1,1,0)&(1,0,1,1)&(2,3,4)\\
2&(0,1,1,1)&(0,1,1,0)&(4,2,3)\\
3&(0,1,1,1)&(0,1,1,2)&(4,4,3)\\
4&(0,1,1,1)&(0,1,2,1)&(4,3,4)\\
5&(0,1,1,1)&(1,0,2,1)&(4,3,2)\\
6&(0,1,2,1)&(0,1,1,1)&(3,4,4)\\
7&(0,1,2,1)&(1,0,0,2)&(3,2,4)\\
8&(0,1,2,1)&(1,1,1,0)&(3,4,2)
\end{array}
 \tag{C.40.5}
\]
There are \(216\) generating projective pairs, and these
are their nine conjugation orbits of size \(24\), as in
C.39. Their passports are pairwise different.

Here is the complete finite verification, with its
necessary-condition direction explicit. For each of
the six choices (C.40.4), consider values
\(F_j,U_j\in L\). Give them the \(S_3\) images prescribed
by (C.40.3) and its transported forms. Retain only the
following finite conditions, all obeyed by the actual
full words \(f,u\).

First, the pair
\[
       \bigl(A_j^{-1},F_j^{-1}B_j^{-1}F_j\bigr)
 \tag{C.40.6}
\]
must be a generating projective pair in orbit \(j\), and
its product must be conjugate in \(Q\) to
\((\bar A_j\bar B_j)^{-1}\). These assertions follow
from the full negative peripheral automorphism and the
distinct passports. Let \(\eta_j(F_j)\) be the conjugation
automorphism of \(L\) transporting the original pair
to (C.40.6) up to central input signs. It is well defined
on \(L\): the projective generating pair has trivial
centralizer, and the remaining central sign acts trivially.
The full involutivity and R1 identities give
\[
             F_j\eta_j(F_j)(F_j)=1,\qquad
             U_j=F_j\eta_j(F_j)(U_j).
 \tag{C.40.7}
\]
Central input changes have no effect on these word
evaluations, even when the normalized action does not
descend to an automorphism of the chosen binary quotient.

Second, impose binary duality for both arrays. For any
generating inputs \(A,B\), this means
\[
                      F(A,B)F(B,A)=1,\qquad
                      U(A,B)U(B,A)=1,
\]
where each value is transported from (C.40.5).
For \(Z=(AB)^{-1}\), impose the actual negative hexagon
\[
                  F(A,B)A^{-1}F(Z,A)Z^{-1}F(B,Z)B^{-1}=1.
 \tag{C.40.8}
\]
Each pair in this equation is still projectively generating,
so every term is among the nine transported values.
Finally retain the already proved C.39 condition that
\((\bar A,\overline{U(A,B)^{-1}BU(A,B)})\) generates \(Q\).

These are constraints on the values of given full words;
no finite-table realization theorem is invoked.
They involve only the \(24\) determinant-one matrices.
Duality and cyclic permutation split the nine indices
into the two components
\[
          \{0,1,2,5,7,8\},\qquad \{3,4,6\}.
\]
Enumerating the indicated matrices, transporting by the
\(48\) invertible matrices, and multiplying in
\(\mathbb F_3\) gives the following complete counts.
The entries in the two middle columns count pairs of
arrays \((F,U)\) on the respective components.
\[
\begin{array}{c|c|c|c}
(k,\ell)&\{0,1,2,5,7,8\}&\{3,4,6\}
       &\text{failures of the }d\text{-third cusp}\\ \hline
(0,0)&64&16&0\\
(0,1)&64&16&0\\
(0,2)&64&16&0\\
(2,0)&64&16&0\\
(2,1)&64&16&0\\
(2,2)&64&16&0
\end{array}
 \tag{C.40.9}
\]
More precisely, each row has four possible \(F\)-arrays
on the first component and two on the second. Every
surviving array satisfies, at every index,
\[
 \overline{A_j^\mu U_j^{-1}B_j^\mu U_j}
       \sim_Q\overline{(A_jB_j)^\mu},
 \qquad
 \mu=\begin{cases}1&k=0,\\5&k=2.\end{cases}
 \tag{C.40.10}
\]
The self-contained certificate
\path{work_CE/code/check_S4_third_cusp.py}
reproduces (C.40.5)--(C.40.10), asserts the complete counts
and zero failures, and retains its exact output alongside
the source. Its enumeration first chooses \(F_j\) satisfying
the local head and norm constraints, imposes transported
duality, then tests (C.40.8), and finally enumerates all
allowed \(U_j\). Thus it covers the entire finite
necessary-condition set. Additional full-word compatibility
could remove entries from that set and cannot create a failure.

The exponent of \(Q=S_4\) is twelve, so the representatives
\(\mu=1,5\) in (C.40.10) suffice for the actual \(d\).
The check is made after projective reduction, and does
not assume that these two integers give \(d\)'s exponent
on the larger binary group \(G\). Since the actual
canonical words obey all the imposed conditions, this
proves (C.40.2) for every epimorphism \(q:F\to S_4\).

\emph{Proper images.}\par
Let \(M=q(F)<S_4\). Under \(S_4\to S_3\), a proper
image has either proper projection, in which case it
lies in \(A_4\) or a Sylow two-subgroup \(D_8\), or
surjective projection, in which case it is an \(S_3\)
point stabilizer by the proof of C.39.
For an \(A_4\) image, C.36 proves the identity for every
unit. In \(D_8\), the commutator subgroup is central,
so \(q(u)\) centralizes the image; its exponent is
four and \(d\equiv1\pmod4\), giving equality.

For an \(S_3\) image generated by two reflections,
the universal dihedral quotient gives exactly
\[
          q(x^d u^{-1}y^d u)
              =(q(x)q(y))^{\,2\kappa+1}
              =q((xy)^d).
\]
For a mixed rotation/reflection generating pair, both
products are reflections and hence are conjugate.
Abelian images have \(q(u)=1\) and give equality.
This proves (C.40.2) for every homomorphism.

\emph{The actual joint finite head.}\par
For every epimorphism to \(S_4\), (C.40.2), the unit
powers on the first two generators, and C.39 show that
the transformed pair is in the same one of the nine
orbits (C.40.5). Thus \(q\theta_d\) has precisely the
same kernel as \(q\) and induces an inner automorphism
of its \(S_4\) target. It follows that \(\theta_d\)
induces an injective endomorphism of the finite group
\(F/J_4\), hence an automorphism.

The basis-power map \(p_d:x\mapsto x^d,y\mapsto y^d\)
is an automorphism of \(F\), with inverse \(p_{d^{-1}}\).
It permutes all epimorphism kernels to \(S_4\), so it
also induces an automorphism of \(F/J_4\).
Since \(\theta_d=\psi p_d\), the canonical \(\psi\)
induces an automorphism of this joint finite quotient
as well. This uses the proved marking and the actual
kernel orbits; it does not deduce joint generation
merely from separate surjectivity tests.
\end{proof}
